% Extracted from Paper_8_Bond_Sphere_Sturmian.tex (Section XV)
% Moved to archive during v1.0.0 cleanup, March 12, 2026
% This section documents a negative result: no D-matrix critical points
% map to R_eq ~ 3.015 bohr for LiH.

\section{Harmonic Phase Locking: A Negative Result}
\label{sec:phase_lock}

\subsection{Hypothesis}

If the equilibrium bond length $R_{\text{eq}}$ is determined by $SO(4)$
representation theory alone, then a critical point (maximum, minimum, or
steepest inflection) of some $D$-matrix element
$D^n_{(l',m'),(l,m)}(\gamma)$ should map to $R \approx 3.015$~bohr
(the experimental LiH equilibrium distance) under a physically motivated
$p_0$ via the inverse stereographic map $R = 1/(p_0 |\tan(\gamma/2)|)$.
Such a correspondence would constitute a closed-form geometric prediction
of $R_{\text{eq}}$ from the $SO(4)$ Wigner $D$-matrix spectrum.

\subsection{Method}

Five $D$-matrix elements were evaluated on a dense grid of $10{,}000$
points in $\gamma \in [0.01, \pi - 0.01]$:
$D^1_{00,00}$, $D^2_{00,00}$, $D^2_{10,10}$, $D^2_{00,10}$, and
$D^3_{00,00}$.
Critical points (zero-crossings of the numerical derivative and second
derivative) were located by interpolation and mapped to physical~$R$ via
$R^* = 1/(p_0 |\tan(\gamma^*/2)|)$ for five physically motivated $p_0$
values: $0.65$ ($Z_{\text{eff}}/n$ for Li $2s$), $0.92$
($\sqrt{2 \times 0.42}$), $1.00$ (atomic H), $1.92$
($\sqrt{2|E_{\text{Li},2s}|}$), and $2.00$ (converged dual-$p_0$ from
v0.9.34).

\subsection{Result}

\begin{enumerate}

\item $D^1_{00,00} = 1$ (constant): no critical points.

\item $D^2_{10,10} = 1$ (constant, Theorem~\ref{thm:selection}):
  no critical points.

\item $D^2_{00,10} = 0$ (constant, Theorem~\ref{thm:selection}):
  no critical points.

\item $D^2_{00,00} = \cos\gamma$: single inflection at
  $\gamma = \pi/2$, mapping to $R^* = 1/p_0$.  At $p_0 = 1.0$:
  $R^* = 1.0$~bohr; at $p_0 = 0.65$: $R^* = 1.54$~bohr.  No value
  maps near 3.015~bohr.

\item $D^3_{00,00}$: minimum at $\gamma = \pi/2$ ($R^* = 1/p_0$,
  same as above); inflections at $\gamma \approx 0.79$ and $2.36$.
  The nearest $R^*$ to 3.015~bohr occurs at $p_0 = 0.65$,
  $R^* = 3.71$~bohr (23\% error).

\end{enumerate}

No strong hits ($R^* \in [2.9, 3.1]$~bohr) and no near hits
($R^* \in [2.8, 3.2]$~bohr) were found for any element at any $p_0$.

At $R_{\text{eq}} = 3.015$~bohr, the derivatives $dD/d\gamma$ are large
and non-zero for all non-trivial elements: $dD^2_{00,00}/d\gamma \in
[-0.81, -0.32]$ and $dD^3_{00,00}/d\gamma \in [-1.31, -0.81]$
across the tested $p_0$ values.  The $D$-matrices sit on smooth,
steep slopes at $\gamma(R_{\text{eq}})$.

\subsection{Conclusion}

The equilibrium bond length $R_{\text{eq}}$ is \emph{not} a geometric
fixed point of the $SO(4)$ representation.  It is an emergent variational
balance---the point where the attractive cross-nuclear potential, the
kinetic repulsion, electron--electron repulsion, and nuclear--nuclear
repulsion reach a global energy minimum.  The $D$-matrices provide the
coupling \emph{strengths} at each $R$, but the equilibrium is a
dynamical property of the full many-body Hamiltonian, not a
representation-theoretic resonance.

This negative result is consistent with the Sturmian structural theorem
(Theorem~\ref{thm:structural}): since the Sturmian eigenvalues are
$\gamma$-independent, the bond sphere geometry alone cannot predict
$R_{\text{eq}}$.  The equilibrium arises from the atom-centered LCAO
architecture, where the distinct atomic length scales
$p_0(A) \neq p_0(B)$ introduce $R$-dependent competition between
attraction and repulsion that the single-$S^3$ Sturmian basis cannot
access.
